% Generated by results/thesis_tables.py - do not edit by hand.
\begin{figure}[H]
\centering
\begin{subfigure}[b]{0.24\textwidth}
\centering
\includegraphics[width=\linewidth]{images/chapter5/confusion_sum.png}
\caption{Πλήρες σύστημα}
\end{subfigure}
\hfill
\begin{subfigure}[b]{0.24\textwidth}
\centering
\includegraphics[width=\linewidth]{images/chapter5/confusion_samples.png}
\caption{Μόνο δείγματα}
\end{subfigure}
\hfill
\begin{subfigure}[b]{0.24\textwidth}
\centering
\includegraphics[width=\linewidth]{images/chapter5/confusion_policy.png}
\caption{Μόνο πολιτικές}
\end{subfigure}
\hfill
\begin{subfigure}[b]{0.24\textwidth}
\centering
\includegraphics[width=\linewidth]{images/chapter5/confusion_llm.png}
\caption{Χωρίς RAG}
\end{subfigure}
\caption{Confusion matrices των τεσσάρων διαμορφώσεων στο στάδιο του LLM. Γραμμές: πραγματική κλάση· στήλες: απόφαση του συστήματος (MAL = \texttt{MALICIOUS}, BEN = \texttt{BENIGN})· σε κάθε κελί το πλήθος πακέτων και το ποσοστό ανά πραγματική κλάση, με την ένταση του χρώματος ανάλογη του ποσοστού.}
\label{fig:confusion_grid}
\end{figure}
